@@M0@@
where the @@M1@@ are the classes of invertible sheaves: to show this, the argument of 3.1 reduces us to the case of a projective bundle over @@M2@@. This last case is an immediate consequence of 1.13. The element @@M3@@ under consideration belongs to @@M4@@ by what precedes, and consequently @@M5@@. Now @@M6@@; thus, by 6.6, @@M7@@. Since @@M8@@ belongs to a basis of @@M9@@ over @@M10@@, this implies @@M11@@. Since @@M12@@ is injective (3.1), we finally have @@M13@@.

\sgaStableTarget{sga6.e06.problem.6-10}\textbf{Problem 6.10.} Under the hypotheses of 6.9, is it true that, for every @@M14@@, we have: @@M15@@?

\sgaStableTarget{sga6.e06.bibliography-for-expose-vi}\subsection*{BIBLIOGRAPHY}
\begin{enumerate}[label={\textup{[\arabic*]}}]
\item Bass H., K-Theory and stable algebra, Publ. Math. de l'IHES, n°22.
\item Chevalley C., Séminaire 1958 : Anneaux de Chow et applications.
\item Hakim M., Topos localement annelés et schémas relatifs, thesis (in preparation).
\item Hartshorne R., Residues and duality, Lecture Notes in Mathematics n°20, Springer-Verlag.
\item Serre J.-P., Le groupe de Grothendieck des schémas en groupes réductifs, Publications Mathématiques de l'IHES, n°34.
\end{enumerate}

% Native Expose VII continuation definitions, retained at their boundary.
\providecommand{\Kosz}{\mathrm K}
\providecommand{\Proj}{\operatorname{Proj}}
\providecommand{\Sym}{\operatorname{Sym}}
\providecommand{\gr}{\operatorname{gr}}
\providecommand{\Tor}{\operatorname{Tor}}
\providecommand{\Spec}{\operatorname{Spec}}
\providecommand{\Imf}{\operatorname{Im}}
\providecommand{\Ker}{\operatorname{Ker}}
\providecommand{\Id}{\operatorname{Id}}
\providecommand{\rad}{\operatorname{rad}}
\providecommand{\cO}{\mathcal O}
\providecommand{\cS}{\mathcal S}
\providecommand{\cJ}{\mathcal J}
\providecommand{\cN}{\mathcal N}
\providecommand{\bfP}{\mathbf P}
\providecommand{\bfA}{\mathbf A}

% Expose VII additions
\providecommand{\Kosz}{\mathrm K}
\providecommand{\Proj}{\operatorname{Proj}}
\providecommand{\Sym}{\operatorname{Sym}}
\providecommand{\gr}{\operatorname{gr}}
\providecommand{\Tor}{\operatorname{Tor}}
\providecommand{\Spec}{\operatorname{Spec}}
\providecommand{\Imf}{\operatorname{Im}}
\providecommand{\Ker}{\operatorname{Ker}}
\providecommand{\Id}{\operatorname{Id}}
\providecommand{\rad}{\operatorname{rad}}
\providecommand{\cO}{\mathcal O}
\providecommand{\cS}{\mathcal S}
\providecommand{\cJ}{\mathcal J}
\providecommand{\cN}{\mathcal N}
\providecommand{\bfP}{\mathbf P}
\providecommand{\bfA}{\mathbf A}

